% p.75 — الباب الثاني والثلاثين في تمرين الصرف والتركيب والاعراب في الجمل
% Layout: \bab heading across the page, then twocol: two columns of continuous, fully
% vocalised Arabic prose (practice sentences run together without punctuation), one
% paragraph per column; rightcol ends «نَنْتَظِرُ خَلِيْلًا», leftcol starts «كُنَّا مُطْمَئِنِّيْنَ».
% TODO: قَبِلَ printed without shadda (expected قَبَّلَ); اَتَغَدّٰى smudged (letter after غ could be read ش).
\mspage{75}
\bab{الباب الثاني والثلاثين في تمرين الصرف والتركيب والاعراب في الجمل}
\begin{twocol}
\begin{rightcol}
اِذْهَبْ مَعَ وَالِدِيْ اِلَى الْمَسْجِدِ اِسْتَيْقِظْ عَلٰى اسْمِ اللهِ مُبَكِّرًا اَصِلُ اِلَى الْمَدْرَسَةِ فِي الْمِيْعَادِ لَا اَتَأَخَّرُ لَا اَقْرَأُ بَعْدَ صَلَاةِ الْعَصْرِ اِلَى الْمَغْرِبِ اَقْضِي الْيَوْمَ فِيْ مُطَالَعَةِ كِتَابٍ وَمُحَادَثَةٍ مَعَ اَبِيْ وَاُمِّيْ وَاِخْوَتِيْ اَوْ فِيْ زِيَارَةِ قَرِيْبٍ اَوْ عِيَادَةِ مَرِيْضٍ مُرُوْا اَوْلَادَكُمْ بِالصَّلٰوةِ وَاضْرِبُوْهُمْ فَاِذَا تَرَكْتَ صَلٰوةً ضَرَبْتُكَ كُنْتُ تَعَلَّمْتُ كَثِيْرًا مِنَ الْاَدْعِيَةِ وُلِدْتُ صَغِيْرًا ضَعِيْفًا لَا اَقْدِرُ عَلٰى عَمَلٍ وَلَا آكُلُ بِنَفْسِيْ هَلْ اَغْضَبَكَ اَحَدٌ لِمَاذَا لَا تَتَكَلَّمُ وَلَمَّا دَخَلْتُ فِي السَّنَةِ الْخَامِسَةِ كَانَ النَّاسُ يُحِبُّوْنَنِيْ وَيُقَرِّبُوْنَنِيْ اِلَيْهِمْ لِاَنِّيْ وَلَدٌ صَالِحٌ اَنَا اَخْدِمُ اَبِيْ وَاُمِّيْ بِالْمَالِ وَالْبَدَنِ خَرَجْتُ يَوْمَ عُطْلَةٍ مَعَ \unclear{صَبِيَّانِ} لَمْ نَزَلْ نَمْشِيْ فِي الْحَرِّ وَالشَّمْسِ مَا وَجَدْنَا طَعَامًا وَلَا مَاءً جَلَسْنَا فِيْ ظِلِّ شَجَرَةٍ نَنْتَظِرُ خَلِيْلًا
\par
\end{rightcol}
\begin{leftcol}
كُنَّا مُطْمَئِنِّيْنَ اِذْ خَرَجَتْ بَقَرَةٌ وَكَانَ عِنْدِيْ سِكِّيْنٌ صَغِيْرٌ حَادٌّ اَنَا اُرِيْدُ اَنْ اَصِيْدَ اَعْطِنِيْ بُنْدُقِيَّتِيْ فَلَمَّا اَذَّنَ اَفْطَرْتُ بِتَمْرَةٍ وَمَاءٍ \unclear{قَبِلَ} اَبِيْ وَرَضِيَتْ اُمِّيْ وَدَعَتْ اَصْدِقَائِيْ قَدْ رَاَيْتُ مَحْمُوْدًا لَمَّا صَامَ لَبِسَ لِبَاسًا جَدِيْدًا وَصُنِعَتْ لَهُ اَطْعِمَةٌ لَذِيْذَةٌ وَخَطَبَ الْاُسْتَاذُ فِي الْاَخِيْرِ اَنَّ النَّبِيَّ صَلَّى اللهُ عَلَيْهِ وَسَلَّمَ كَانَ يُسَابِقُ وَقَدْ كَانَ سَيِّدُنَا عُثْمَانُ غَنِيًّا وَسَيِّدُنَا عَبْدُ الرَّحْمٰنِ بْنُ عَوْفٍ غَنِيًّا وَقَالَ النَّبِيُّ صَلَّى اللهُ عَلَيْهِ وَسَلَّمَ اَلْمَالُ الصَّالِحُ خَيْرٌ لِلرَّجُلِ الصَّالِحِ جَاءَتْ مَرْكَبَتَانِ فَرَكِبْنَاهُمَا وَسَلَّمْتُ عَلٰى اَبِيْ وَصَلْنَا اِلَى الْمَحَطَّةِ وَاشْتَرَيْنَا التَّذْكِرَةَ كُنْتُ اَخْرُجُ اَحْيَانًا مَعَ الْمُسْلِمِيْنَ اِلَى الْجِهَادِ لَنْ تَجِدَ يَا اَخِيْ اَرْخَصَ مِنْ هٰذَا فِي السُّوْقِ اُصَدِّقُكَ لِاَنَّكَ مُسْلِمٌ وَالْمُسْلِمُ لَا يَكْذِبُ اَلَا تَرٰى كَيْفَ بَسَطُوْا اَنْوَاعًا مِنَ الثِّيَابِ \unclear{اَتَغَدّٰى} مَعَ وَالِدِيْ وَاَحْفَظُ دُرُوْسِيْ وَاُطَالِعُ خَرَجْتُ مَرَّةً اِلٰى مُبَارَاةٍ وَكَانَ الزِّحَامُ شَدِيْدًا اِذَا صِرْتَ عَالِمًا فَمَاذَا تَعْمَلُ لِلْمُسْلِمِيْنَ
\par
\end{leftcol}
\end{twocol}
